\section{Ejercicio integrador: del VDP Report a la divulgación y el aprendizaje}

Para cerrar, un tabletop enlaza los mecanismos de esta clase en una ruta que puede ensayarse. Supongamos que un investigador envía por el VDP un reporte de alto impacto; a primera vista parece solo una vulnerabilidad, pero durante la verificación se descubre el uso real de una credential y un posible acceso a datos. El objetivo del ejercicio no es adivinar los motivos del atacante, sino observar si el equipo puede escalar a tiempo, preservar la evidencia, proteger a los usuarios, completar el análisis de materiality y mantener una comunicación coherente.

\subsection{Tabletop en cinco etapas}

En cada etapa solo se entrega a los participantes una parte de los hechos, y se les pide registrar sus supuestos, el owner y el siguiente paso. El Facilitator debe introducir incertidumbre a propósito: registros faltantes, dependencias de terceros, un investigador que exige respuesta rápida, un área de negocio que quiere mantener el servicio y medios de comunicación que empiezan a hacer preguntas. Una buena respuesta no es «cerrar todo de inmediato» ni «esperar a que termine la investigación», sino una que explique el beneficio, el costo, la evidencia requerida y quién aprueba cada acción.

\lecturefigure{16-tabletop.png}{El tabletop enlaza VDP, incident, materiality, notification y learning en un ejercicio de extremo a extremo.}{Diseño integrador de esta clase; concepto redibujado.}

\begin{knowledgebox}{Lectura de la figura: cada etapa debe entregar un resultado con capacidad de auditoría}
El VDP report entrega el case ID, el scope y el acknowledgement; el incident entrega la severity, el IC, la containment y el evidence plan; la governance entrega el materiality owner, la RACI y el decision log; notify/learn entrega mensajes coherentes, el deadline, la recovery y los remediation owners. El éxito del ejercicio no se mide por «no haber discutido», sino por si la discusión giró en torno a los mismos hechos, si se escaló antes del deadline y si la decisión la tomó el rol correcto.
\end{knowledgebox}

\begin{importantbox}{Tabletop injects}
\begin{enumerate}
  \item El reporte proviene de un out-of-country researcher, describe una vulnerabilidad in-scope, pero adjunta registros reales de usuarios;
  \item los internal logs muestran que la credential ya se había usado, en una fecha anterior al reporte;
  \item un SaaS de terceros también podría estar afectado, con una ventana contractual de notificación distinta;
  \item la corrección requiere una breve interrupción del servicio, y el responsable de negocio pide posponerla;
  \item el counsel solicita un análisis de materialidad, y communications recibe preguntas de un periodista;
  \item nueva evidencia contradice la cifra inicial de registros, y el equipo debe corregir, no sobrescribir, la decisión anterior.
\end{enumerate}
\end{importantbox}

\begin{warningbox}{El ejercicio no debe buscar una «respuesta correcta única»}
Las obligaciones reales dependen de los hechos y de la jurisdicción. El tabletop debe poner a prueba las interfaces: si se identifica el trigger de escalamiento, si se preserva la evidence clave antes y después de la containment, si el GC, el CEO o el committee se incorporan cuando hace falta, y si se pueden explicar las decisiones ante los usuarios y el público. Si el ejercicio solo evalúa la memorización de políticas, no podrá revelar los puntos de ruptura de la organización.
\end{warningbox}

\subsection{Resumen de la sección}

El ejercicio de extremo a extremo convierte las políticas, de documentos, en músculo organizacional. Solo cuando VDP, incident command, evidence, materiality, executive decision y learning colaboran en un mismo escenario puede el equipo descubrir los vacíos de responsabilidad que existen en la realidad.
